\documentclass[12pt,a4paper]{article}
\usepackage[utf8]{inputenc}
\usepackage[spanish]{babel}
\usepackage{geometry}
\usepackage{listings}
\usepackage{xcolor}
\usepackage{parskip}
\usepackage{amsmath}
\usepackage{hyperref}

\geometry{margin=2.5cm}

\lstset{
  basicstyle=\ttfamily\small,
  keywordstyle=\color{blue}\bfseries,
  commentstyle=\color{gray},
  breaklines=true,
  frame=single,
  language=Java,
  numbers=left,
  numberstyle=\tiny\color{gray},
  literate={á}{{\'a}}1 {é}{{\'e}}1 {í}{{\'\i}}1 {ó}{{\'o}}1 {ú}{{\'u}}1 {ñ}{{\~n}}1 {Á}{{\'A}}1 {É}{{\'E}}1 {Ó}{{\'O}}1
}

\title{\textbf{Misioneros y Caníbales}\\
  \large BFS, DFS y DFS Recursivo --- Documentación Técnica}
\author{Inteligencia Artificial}
\date{}

\begin{document}
\maketitle
\tableofcontents
\newpage

\section{Descripción del Problema}
El problema de \textbf{Misioneros y Caníbales} es un rompecabezas clásico de IA. Tres misioneros y tres caníbales deben cruzar un río en un bote con capacidad máxima de 2 personas, cumpliendo la restricción de que los caníbales nunca superen en número a los misioneros en ninguna orilla.

\subsection{Representación del estado}
Cada estado se codifica como un vector de 5 enteros:
\[ [m_{izq},\ c_{izq},\ bote,\ m_{der},\ c_{der}] \]

\begin{itemize}
  \item \textbf{Estado inicial}: $[3, 3, 0, 0, 0]$
  \item \textbf{Estado meta}: $[0, 0, 1, 3, 3]$
  \item \textbf{bote}: 0 = orilla izquierda, 1 = orilla derecha
\end{itemize}

\subsection{Restricciones de validez}
Un estado es válido si:
\begin{enumerate}
  \item Todos los valores son $\geq 0$ y $\leq 3$.
  \item Si $m_{izq} > 0$, entonces $m_{izq} \geq c_{izq}$.
  \item Si $m_{der} > 0$, entonces $m_{der} \geq c_{der}$.
\end{enumerate}

\section{Movimientos Posibles}
Desde cada estado se generan hasta 5 movimientos posibles (lo que cabe en el bote):
\begin{center}
\begin{tabular}{|c|c|}
\hline
\textbf{Misioneros} & \textbf{Caníbales} \\
\hline
1 & 0 \\
2 & 0 \\
0 & 1 \\
0 & 2 \\
1 & 1 \\
\hline
\end{tabular}
\end{center}

\section{BFS --- Búsqueda por Amplitud}

BFS usa una \textbf{cola FIFO} para explorar los estados nivel por nivel, garantizando la solución óptima (menor número de pasos).

\begin{lstlisting}
Nodo_BFS resolver() {
    cola.encolar(estadoInicial);
    visitados.add(estadoInicial);

    while (!cola.estaVacia()) {
        Nodo_BFS actual = cola.desencolar();
        if (esSolucion(actual)) return actual;

        for (Nodo_BFS hijo : generarHijos(actual)) {
            if (!estaVisitado(hijo)) {
                visitados.add(hijo);
                cola.encolar(hijo);
            }
        }
    }
    return null;
}
\end{lstlisting}

\textbf{Resultado}: solución óptima en \textbf{11 pasos}, 15 nodos generados.

\section{DFS --- Búsqueda por Profundidad con Pila}

DFS usa una \textbf{pila LIFO} para explorar el camino más profundo primero. No garantiza la solución óptima pero usa menos memoria que BFS.

\begin{lstlisting}
Nodo_DFS resolver() {
    pila.apilar(estadoInicial);
    visitados.add(estadoInicial);

    while (!pila.estaVacia()) {
        Nodo_DFS actual = pila.desapilar();
        if (esSolucion(actual)) return actual;

        for (Nodo_DFS hijo : generarHijos(actual)) {
            if (!estaVisitado(hijo)) {
                visitados.add(hijo);
                pila.apilar(hijo);
            }
        }
    }
    return null;
}
\end{lstlisting}

\section{DFS Recursivo}

Implementación recursiva del DFS. Más elegante en código pero con riesgo de \texttt{StackOverflowError} en problemas más profundos.

\begin{lstlisting}
Nodo_DFS resolverRecursivo(Nodo_DFS nodo,
                            Set<String> visitados) {
    if (esSolucion(nodo)) return nodo;

    for (Nodo_DFS hijo : generarHijos(nodo)) {
        String clave = estadoAClave(hijo);
        if (!visitados.contains(clave)) {
            visitados.add(clave);
            Nodo_DFS resultado =
                resolverRecursivo(hijo, visitados);
            if (resultado != null) return resultado;
        }
    }
    return null;
}
\end{lstlisting}

\section{Reconstrucción del Camino}

Cada nodo guarda una referencia al nodo padre. Para reconstruir el camino desde la solución hasta el estado inicial se recorre la cadena de padres:

\begin{lstlisting}
List<int[]> reconstruirCamino(Nodo_BFS solucion) {
    List<int[]> camino = new ArrayList<>();
    Nodo_BFS nodo = solucion;
    while (nodo != null) {
        camino.add(0, nodo.getEstado());
        nodo = nodo.getPadre();
    }
    return camino;
}
\end{lstlisting}

\section{Interfaz Gráfica}

La clase \texttt{InterfazCyM} implementa:
\begin{itemize}
  \item Representación visual de ambas orillas con personajes dibujados en \texttt{paintComponent()}.
  \item Botones para mover 1M, 2M, 1C, 2C ó 1M+1C manualmente.
  \item Botón \textbf{``Resolver (BFS)''} que ejecuta BFS y anima el camino solución con un \texttt{javax.swing.Timer} de 700ms entre pasos.
\end{itemize}

\end{document}
